\section{Introdução}

A visão computacional tornou-se um dos pilares fundamentais da tecnologia moderna, impulsionando inovações que variam desde sistemas de condução autônoma e diagnósticos médicos por imagem até autenticação biométrica em dispositivos móveis \citep{Szeliski2022}. No centro dessa revolução encontra-se a biblioteca OpenCV (\textit{Open Source Computer Vision Library}), que se estabeleceu como o padrão \textit{de facto} na indústria e na academia devido à sua robustez e vasta coleção de algoritmos otimizados \citep{bradski2000opencv}.

Embora o núcleo do OpenCV seja desenvolvido em C++ para garantir máxima eficiência computacional, a biblioteca oferece interfaces (\textit{bindings}) para diversas outras linguagens, como Python e Java. Essa flexibilidade impõe aos engenheiros e pesquisadores um dilema arquitetural documentado: a necessidade de balancear a eficiência bruta e o controle de memória do C++ com a produtividade e a abstração de alto nível oferecidas por linguagens gerenciadas \citep{Prechelt, 2000}.



A popularização de linguagens como Python, impulsionada pelo ecossistema de Ciência de Dados, e a onipresença do Java em sistemas corporativos, trazem à tona o debate sobre o "custo da abstração". A literatura indica que camadas adicionais de interpretação, o gerenciamento automático de memória (\textit{Garbage Collection}) e os mecanismos de interoperação introduzem sobrecargas computacionais (overhead) que podem comprometer o desempenho em tempo real de sistemas críticos \citep{Pereira et al., 2017; Nanz & Furia, 2015}.
Este trabalho tem como objetivo principal quantificar e analisar esse impacto. Através de um estudo comparativo, avalia-se o desempenho do algoritmo de detecção de placas de carro, utilizando modelos \textit{Haar Cascade}, implementado em C++, Java e Python. A análise não se restringe apenas à latência média, mas investiga também qualitativamente a usabilidade, disponibilidade de recursos e a curva de aprendizado de cada ambiente.

As contribuições deste artigo visam fornecer diretrizes empíricas para auxiliar arquitetos de \textit{software} na escolha da pilha tecnológica mais adequada, balanceando os requisitos não-funcionais de desempenho com as necessidades de negócio e produtividade.

Este trabalho está estruturado da seguinte forma: a seção 2 apresenta uma revisão bibliográfica sobre as linguagens de programação e biblioteca estudadas. A seção 3 detalha a metodologia adotada para a implementação e avaliação dos algoritmos. A seção 4 apresenta os resultados obtidos. A seção 5 discute os resultados obtidos, e a seção 6 conclui o artigo, destacando as principais descobertas e sugerindo direções para trabalhos futuros.